\documentclass[11pt]{article}

\usepackage[T1]{fontenc}
\usepackage[utf8]{inputenc}
\usepackage{lmodern}
\usepackage{amsmath,amssymb,amsthm,mathtools}
\usepackage[margin=1in]{geometry}
\usepackage[colorlinks=true,linkcolor=blue,urlcolor=blue]{hyperref}
\usepackage{microtype}

\newtheorem{theorem}{Theorem}
\newtheorem{proposition}{Proposition}
\newtheorem{remark}{Remark}
\newcommand{\dd}{\,\mathrm{d}}
\newcommand{\avg}[1]{\langle #1\rangle_w}
\newcommand{\Var}{\operatorname{Var}}

\title{Spectral RMS Is Not Generally Equal to Averaged Curvature\\
\large An Exact Spectral--Phase Identity and a Closed-Curve Counterexample}
\author{Richard J. Reyes}
\date{October 26, 2026}

\begin{document}
\maketitle

\begin{abstract}
For a smooth, nonvanishing eigenmode of the periodic operator
$-\partial_s^2$, we establish an exact decomposition of its eigenvalue
into squared weighted mean phase gradient, weighted phase variance, and
normalized amplitude-gradient energy.  The decomposition yields sharp
equality conditions and does not require a nonnegative phase gradient.
An explicit eigenmode with $\sqrt{\lambda}=1$ and weighted mean phase
gradient $4/5$ can moreover be realized as the unit tangent of a smooth
closed planar curve, where the phase gradient equals its curvature.
The conclusion applies to the declared free periodic operator, without
assuming a derived physical confinement or mass operator.
\end{abstract}

\section{Setting and exact identity}

Let $S_L^1=\mathbb{R}/L\mathbb{Z}$ denote a circle of length $L>0$.
Consider the self-adjoint free operator
\begin{equation}
    H=-\frac{\mathrm{d}^2}{\mathrm{d}s^2},
    \qquad D(H)=H^2_{\mathrm{per}}(S_L^1;\mathbb{C}).
\end{equation}
Suppose that $\chi$ is a smooth, nowhere-vanishing periodic eigenfunction:
\begin{equation}\label{eq:eigenproblem}
    -\chi''(s)=\lambda\chi(s),\qquad \chi(s)\ne0.
\end{equation}
Because $\chi$ has no zeros, its polar form may be written on the
universal cover as
\begin{equation}\label{eq:polar}
    \chi(s)=R(s)e^{i\phi(s)},\qquad R(s)>0,\qquad
    \sigma(s)=\phi'(s)\in\mathbb{R}.
\end{equation}
Although $\phi$ need not be periodic, $\phi'$ is periodic, and
$\phi(s+L)-\phi(s)=2\pi n$ for some $n\in\mathbb{Z}$.
Set
\begin{equation}\label{eq:weights}
    N=\int_0^L R(s)^2\dd s>0,\qquad
    w(s)=\frac{R(s)^2}{N},\qquad
    \avg{f}=\int_0^L f(s)w(s)\dd s,
\end{equation}
and define
\begin{equation}
    \Var_w(\sigma)=\avg{\sigma^2}-\avg{\sigma}^{\,2}
    =\int_0^L w(s)\bigl(\sigma(s)-\avg{\sigma}\bigr)^2\dd s.
\end{equation}

\begin{theorem}[Sharp spectral--phase decomposition]\label{thm:main}
Under the hypotheses above,
\begin{equation}\label{eq:identity}
    \boxed{
      \lambda-\avg{\sigma}^{\,2}
      =\frac{\displaystyle\int_0^L |R'(s)|^2\dd s}
             {\displaystyle\int_0^L R(s)^2\dd s}
       +\Var_w(\sigma)\ge0.
    }
\end{equation}
Equality in \eqref{eq:identity} holds if and only if both $R$ and
$\sigma$ are constant.  In particular,
\begin{equation}\label{eq:rootbound}
    \sqrt{\lambda}\ge\bigl|\avg{\sigma}\bigr|.
\end{equation}
No sign restriction on $\sigma$ is needed.
\end{theorem}

\begin{proof}
Multiply \eqref{eq:eigenproblem} by $\overline{\chi}$ and integrate:
\begin{equation}
    \lambda\int_0^L |\chi|^2\dd s
      =-\int_0^L\overline{\chi}\,\chi''\dd s.
\end{equation}
Periodic integration by parts, using periodicity of both $\chi$ and
$\chi'$, eliminates the boundary term and yields
\begin{equation}\label{eq:dirichlet}
    \lambda N=\int_0^L|\chi'|^2\dd s.
\end{equation}
From \eqref{eq:polar},
\begin{equation}
    \chi'=e^{i\phi}(R'+iR\sigma),
    \qquad
    |\chi'|^2=|R'|^2+R^2\sigma^2,
\end{equation}
since $R'$ and $\sigma$ are real.  Inserting this expression into
\eqref{eq:dirichlet} and dividing by $N$ gives
\begin{equation}\label{eq:phase_second_moment}
    \lambda
    =\frac{\int_0^L|R'|^2\dd s}{N}+\avg{\sigma^2}.
\end{equation}
By the weighted variance identity,
$\avg{\sigma^2}=\avg{\sigma}^{\,2}+\Var_w(\sigma)$.
Substitution into \eqref{eq:phase_second_moment} proves
\eqref{eq:identity}.  Both terms on its right-hand side are
nonnegative, proving \eqref{eq:rootbound}.

For the equality case, the sum in \eqref{eq:identity} vanishes if and
only if both of its nonnegative terms vanish.  As $R$ is smooth,
$\int_0^L|R'|^2\dd s=0$ implies $R'=0$ everywhere, hence $R$ is
constant.  Likewise, $w>0$ everywhere and smoothness of $\sigma$
show that $\Var_w(\sigma)=0$ implies
$\sigma(s)=\avg{\sigma}$ everywhere.  The converse follows by direct
substitution, completing the proof.
\end{proof}

\begin{remark}[Winding and the sign of the mean]
When equality holds, $\chi(s)=R_0e^{i(\sigma_0s+\phi_0)}$, where
$\sigma_0L=2\pi n$, $n\in\mathbb{Z}$, and
$\lambda=\sigma_0^2$.  Consequently, equality in
$\lambda=\avg{\sigma}^{\,2}$ allows either sign of the constant
phase gradient.  The stronger unsigned statement
$\sqrt{\lambda}=\avg{\sigma}$ requires in addition
$\avg{\sigma}\ge0$; without this assumption the exact equality is
$\sqrt{\lambda}=|\avg{\sigma}|$.
\end{remark}

\section{An explicit strict counterexample}

\begin{proposition}[Nonconstant amplitude defeats mean closure]
\label{prop:example}
For $s\in[0,2\pi]$, let
\begin{equation}\label{eq:example}
    \chi_a(s)=\cos s+ia\sin s,\qquad a>0.
\end{equation}
Then $\chi_a$ is a smooth, nowhere-vanishing periodic eigenfunction
with $\lambda=1$, while its weighted mean phase gradient is
\begin{equation}\label{eq:generalmean}
    \avg{\sigma}=\frac{2a}{1+a^2}.
\end{equation}
Thus $\sqrt{\lambda}\ne\avg{\sigma}$ for every $a>0$ with
$a\ne1$.  In particular, at $a=\tfrac12$,
\begin{equation}\label{eq:fourfifths}
    \boxed{\sqrt{\lambda}=1,\qquad\avg{\sigma}=\frac45.}
\end{equation}
\end{proposition}

\begin{proof}
Direct differentiation yields $-\chi_a''=\chi_a$, hence $\lambda=1$.
Moreover,
\begin{equation}
    R_a^2(s)=|\chi_a(s)|^2=\cos^2s+a^2\sin^2s
    \ge\min\{1,a^2\}>0.
\end{equation}
The normalization is
\begin{equation}
    N_a=\int_0^{2\pi}R_a^2(s)\dd s=\pi(1+a^2).
\end{equation}
For a nonvanishing complex function, its phase derivative is
\begin{equation}
    \sigma_a(s)=\frac{\operatorname{Im}
    (\overline{\chi_a(s)}\chi_a'(s))}{|\chi_a(s)|^2}.
\end{equation}
Because $\operatorname{Im}(\overline{\chi_a}\chi_a')=a$,
we have
\begin{equation}\label{eq:sigmaa}
    \sigma_a(s)=\frac{a}{R_a^2(s)}>0.
\end{equation}
Therefore $R_a^2\sigma_a=a$ pointwise and
\begin{equation}
    \avg{\sigma}
      =\frac{1}{N_a}\int_0^{2\pi}R_a^2\sigma_a\dd s
      =\frac{2\pi a}{\pi(1+a^2)}
      =\frac{2a}{1+a^2}.
\end{equation}
The difference
\begin{equation}
    1-\frac{2a}{1+a^2}=\frac{(a-1)^2}{1+a^2}
\end{equation}
is strictly positive unless $a=1$.  Taking $a=1/2$ gives
\eqref{eq:fourfifths}.
\end{proof}

\paragraph{Exact gap accounting at $a=1/2$.}
For the explicit example $R^2=\cos^2s+\tfrac14\sin^2s$ and
$N=5\pi/4$.  The elementary integral
\begin{equation}
    \int_0^{2\pi}\frac{\dd s}
    {\cos^2s+\frac14\sin^2s}=4\pi
\end{equation}
gives, using $\sigma=(2R^2)^{-1}$,
\begin{equation}
    \avg{\sigma^2}
      =\frac{1}{4N}\int_0^{2\pi}\frac{\dd s}{R^2(s)}
      =\frac45.
\end{equation}
Hence
\begin{equation}
    \Var_w(\sigma)=\frac45-\frac{16}{25}=\frac4{25},
\end{equation}
while \eqref{eq:phase_second_moment} yields
\begin{equation}
    \frac{\int_0^{2\pi}|R'|^2\dd s}{N}
      =1-\frac45=\frac15.
\end{equation}
Thus both contributions to the gap are explicit:
\begin{equation}\label{eq:gap_accounting}
    \boxed{
    1-\left(\frac45\right)^2
       =\underbrace{\frac15}_{\text{amplitude-gradient energy}}
       +\underbrace{\frac4{25}}_{\text{phase variance}}
       =\frac9{25}.}
\end{equation}

\section{Realization by a smooth closed planar curve}

\begin{proposition}[Geometrically locked counterexample]
\label{prop:curve}
For $a>0$, define the complex unit tangent
\begin{equation}
    T_a(s)=\frac{\chi_a(s)}{|\chi_a(s)|}
\end{equation}
and the plane curve, with $\mathbb{C}\cong\mathbb{R}^2$,
\begin{equation}
    \Gamma_a(s)=\int_0^s T_a(t)\dd t.
\end{equation}
Then $\Gamma_a$ is a smooth, closed, unit-speed planar curve of
length $2\pi$.  Its curvature $\kappa_a$ and torsion $\tau_a$ satisfy
\begin{equation}
    \kappa_a(s)=\sigma_a(s)>0,\qquad \tau_a(s)=0,
\end{equation}
and hence the pointwise geometric lock
\begin{equation}
    \sigma_a(s)=\sqrt{\kappa_a(s)^2+\tau_a(s)^2}
\end{equation}
holds.  Nevertheless, at $a=1/2$ the associated free periodic
eigenmode has $\sqrt{\lambda}=1$ but
$\avg{\kappa}=4/5$.
\end{proposition}

\begin{proof}
Smoothness follows from $|\chi_a|>0$, and
$|\Gamma_a'|=|T_a|=1$.  Since
$\chi_a(s+\pi)=-\chi_a(s)$ and $R_a(s+\pi)=R_a(s)$,
we have $T_a(s+\pi)=-T_a(s)$.  It follows that
\begin{equation}
    \Gamma_a(2\pi)-\Gamma_a(0)
    =\int_0^\pi\bigl(T_a(s)+T_a(s+\pi)\bigr)\dd s=0.
\end{equation}
The curve is closed.  Writing $T_a=e^{i\phi_a}$,
\begin{equation}
    T_a'=i\phi_a'T_a=i\sigma_a T_a,
\end{equation}
so its curvature is $\kappa_a=|T_a'|=|\sigma_a|=\sigma_a$,
where positivity follows from \eqref{eq:sigmaa}.  As the curve is
planar, its torsion is zero.  The final assertion follows from
Proposition~\ref{prop:example}.
\end{proof}

\section{Scope and consequences}

The exact relation can be written in spectral-root form as
\begin{equation}\label{eq:rootidentity}
    \boxed{
    \sqrt{\lambda}=
    \left(
      \avg{\sigma}^{\,2}
      +\Var_w(\sigma)
      +\frac{\|R'\|_{L^2}^2}{\|R\|_{L^2}^2}
    \right)^{1/2}.}
\end{equation}
This demonstrates that the spectral scale $\sqrt{\lambda}$
cannot generally be replaced by the weighted mean $\avg{\sigma}$,
even if a pointwise identification of $\sigma$ with geometric
curvature has been independently supplied.

If the separated-field dispersion of an independently specified
product-space model yields a conditional effective mass
$m_{\mathrm{eff}}=(\hbar/c)\sqrt{\lambda}$, then replacing this with
$(\hbar/c)\avg{\sigma}$ demands the equality hypotheses, together
with the appropriate sign convention.  This theorem does not derive
the confinement operator from an underlying physical action, identify
particle modes, or establish an empirical mass spectrum.  Operators
with nonzero potentials, connections, or different boundary conditions
may require a different energy identity; zero-containing eigenmodes
require additional treatment at their nodes.

\paragraph{Source and context.}
See the project note
\href{https://github.com/rickyjreyes/geometry_of_resonance/blob/main/docs/foundational/GEOMETRIC_MASS_CLOSURE.md}
{\texttt{GEOMETRIC\_MASS\_CLOSURE.md}}.

\end{document}